\uuid{sxs2}
\exo7id{4528}
\titre{exo7 4528}
\auteur{quercia}
\organisation{exo7}
\datecreate{2010-03-14}
\isIndication{false}
\isCorrection{true}
\chapitre{Suite et série de fonctions}
\sousChapitre{Autre}
\cpge{mpsi-mp2i.fonctions-valeurs-reelles-complexes.a, pcsi.techniques-calcul-fonctions.a, ptsi.techniques-calcul-fonctions.a, mpsi-mp2i.calcul-algebrique-trigonometrie.d, pcsi.calcul-algebrique-trigonometrie.d, ptsi.calcul-algebrique-trigonometrie.d, bcpst1.fonctions-reelles-usuelles.a, bcpst1.trigonometrie, mp-mpi.suites-series-fonctions.a, pc.suites-series-fonctions.a, psi.suites-series-fonctions.a, bcpst2.series-reelles}
\module{Analyse}
\niveau{L2}
\difficulte{}

\contenu{
\texte{
On pose $f(x) = \sum_{n=0}^\infty \frac{\arccos(\cos nx)}{n!}$.
}
\begin{enumerate}
    \item \question{Montrer que $f$ est définie sur $\R$, continue, paire et $2\pi$-périodique.}
    \item \question{Calculer $f(0)$, $f(\pi)$, $f\left(\frac\pi2\right)$.}
\reponse{
$f(0) = 0$, $f(\pi) = \pi\sh 1$,
             $f\left(\frac\pi2\right) = \frac\pi2(e-\cos 1)$.
}
\end{enumerate}
}
